\documentclass[12pt]{article}
\usepackage[utf8]{inputenc}
\usepackage[T1]{fontenc}
\usepackage{hyperref}
\title{Puinave Language Resources for K-12 Curriculum Development}
\author{Compiled by Claude Sonnet 5.5 (claude-sonnet-5-5) for Andr\'as Kornai}
\date{}
\begin{document}
\maketitle

\section{Introduction}
Puinave is an indigenous language of the Inírida River region (Guain\'ia) in Colombia and of the Orinoco border area in Venezuela. It is generally treated as a language isolate or as unclassified; the older Mak\'u-Puinave grouping rests on scant evidence. Web searches found general profiles and one grammatical description, but no school materials. This report is deliberately short.

\section{Primary Educational Resources}
No primer, textbook or school programme for Puinave turned up in the searches. The language is not known to be taught in schools.

\section{Linguistic and Reference Materials}
\begin{itemize}
\item Wikipedia, ``Puinave language'': general overview, speaker figures and classification notes. \url{https://en.wikipedia.org/wiki/Puinave_language}
\item LINGUIST List 20.814: review (in Spanish) of what it describes as the first grammatical description of Puinave, based on data from villages along the Inírida River. \url{https://dev.linguistlist.org/issues/20/814/}
\item MPI-EVA numeral systems page for Puinave (notes the language is tonal and that Spanish numerals have largely replaced native ones). \url{https://lingweb.eva.mpg.de/channumerals/Puinave.htm}
\end{itemize}
No dedicated Puinave dictionary was identified in the search results.

\section{Religious and Community Resources}
None found in the searches.

\section{Digital and Audio Resources}
\begin{itemize}
\item UNESCO Atlas of World Languages, Puinave in Colombia: language profile page. \url{https://ar.wal.unesco.org/en/countries/colombia/languages/puinave}
\item Endangered Languages Project entry for Puinave. \url{https://www.endangeredlanguages.com/lang/2767?hl=ru}
\end{itemize}
No audio or video materials were identified.

\section{Curriculum Development Context}
Documentation is thin: one grammatical description and general profiles. A curriculum would first need community consultation, a practical orthography, and lexical documentation. The grammar above is the most useful starting reference.
\end{document}